\chapter{Role du projet Promptly}
\label{ch:role}

\section{Contexte}

\textbf{Promptly} est une plateforme web d'apprentissage dediee a la programmation
(Java, Spring Boot, Angular, algorithmique). Elle s'adresse aux etudiants, aux
enseignants et a un administrateur.

\section{Fonctionnalites principales}

\begin{itemize}
  \item \textbf{Cours et chapitres} : contenus, PDF, resume IA, flashcards ;
  \item \textbf{Quiz} : QCM, scores, detection GUESS/KNOW (ML) ;
  \item \textbf{Forum} : discussions, auto-correction (TextGears), censure (Sightengine) ;
  \item \textbf{Reclamations} : depot etudiant, labels ML + emotion (admin) ;
  \item \textbf{Recrutement} : offres (admin), candidature + CV, matching competences (Gemini) ;
  \item \textbf{Assistant IA} : chatbot general (Ollama).
\end{itemize}

\section{Pile technique}

\begin{table}[H]
\centering
\begin{tabular}{@{}ll@{}}
\toprule
Couche & Technologie \\
\midrule
Frontend & Angular \\
Backend & Spring Boot 3 \\
Base de donnees & MySQL \\
LLM local & Ollama \\
ML locaux & Python (FastAPI) --- ports 8090 / 8091 \\
\bottomrule
\end{tabular}
\caption{Pile technique de Promptly}
\end{table}

\section{Roles utilisateurs}

\begin{itemize}
  \item \textbf{ETUDIANT} : cours, quiz, forum, reclamations, candidatures, assistant ;
  \item \textbf{ENSEIGNANT} : gestion des cours / quiz (+ meme acces apprenant) ;
  \item \textbf{ADMIN} : utilisateurs, traitement reclamations, offres de recrutement, vues ML.
\end{itemize}
